\begin{abstract}
The United States incarcerates approximately 2.1 million people. Under Census Bureau practice, these individuals are counted as residents of the county containing the prison or jail, not the county of their last home address. Because mass incarceration is spatially concentrated in rural counties that receive prisoners disproportionately from urban areas, this counting convention systematically inflates the census population---and therefore the political representation---of rural legislative districts at the expense of urban communities from which prisoners originate. We call this phenomenon \emph{prison gerrymandering}. Using 2020 Census data matched against Department of Justice facility rosters and Prison Policy Initiative residence records, we quantify the rural inflation effect in the 50 largest prison-dense counties, finding population inflation ranging from 8\% to 22\% relative to the prison-free counterfactual. We present three state case studies: New York (which enacted residence-address correction in 2010 and has used it in both the 2010 and 2020 redistricting cycles), Pennsylvania (which has rejected correction legislation), and Texas (which has rejected correction legislation and hosts the largest prison population of any state). We then analyze how BISECT's algorithmic redistricting pipeline would behave under a \texttt{--balance-metric prison\_adjusted} mode, where prisoner populations are reallocated to their last-known home census tract using Prison Policy Initiative's methodology. This mode is identified as a planned preprocessing extension rather than a currently implemented flag. Our results demonstrate that the choice to count prisoners at their facility location is not neutral: it transfers political power from Black and Hispanic urban communities to white rural communities in a manner consistent with the definition of vote dilution under \emph{Thornburg v. Gingles}.
\end{abstract}
